\uuid{38zW}
\titre{Aire/volume}
\auteur{bodin ; hanani ; mzari}
\organisation{exo7}
\licence{CC BY-NC-SA 4.0}
\datecreate{2020-01-20}
\niveau{L1}
\module{Géométrie}
\chapitre{Géométrie du plan}
\sousChapitre{Distance, norme, produit scalaire}
\difficulte{2}
\theme{aire, volume}
\exo7qcmid{lille-1.geometrie-plane.020}

\contenu{
\begin{question}
\qclassification{140.01,140.02}
\qtags{motcle=aire/volume}
Dans le plan muni d'un repère orthonormé $(O,\vec{i},\vec{j})$, on considère les points $\displaystyle A(a,b)$ et $\displaystyle B(1,1)$. Comment choisir les réels $a$ et $b$ pour que l'aire du triangle de sommets $O,A,B$ soit égale à $1$ ?
\begin{answers}
    \good{$\displaystyle a=2$ et $\displaystyle b=0$}
    \good{$\displaystyle a=2+b$ et $\displaystyle b\in \Rr$}
    \bad{$\displaystyle a=1$ et $\displaystyle b=0$}
    \bad{$\displaystyle a=0$ et $\displaystyle b=1$}
\end{answers}
\begin{explanations}
On doit avoir $\displaystyle 2\mbox{Aire}(OAB)=\left|\mbox{det}\left(\overrightarrow{OA},\overrightarrow{AB}\right)\right|=2$. Ceci donne ($a=2+b$ et $b\in \Rr$).
\end{explanations}
\end{question}
}
